\section{Competing Computations and Registered Count States}
\label{sec:representations}

\subsection{Two Direct hypotheses and one sequential CoT program}

We use minimal equations to specify what each experiment must distinguish,
without assuming how a transformer implements the operation. For Direct,
\begin{equation}
  w_i^D=\mathcal R_D(q,R_i),\qquad
  Z_D=\mathcal G_D(\{w_i^D\}_{i=1}^{M}),\qquad
  \widehat N_D=\mathcal Q_D(Z_D).
\label{eq:direct-program}
\end{equation}
\(\mathcal R_D\) routes record information, \(\mathcal G_D\) constructs a
terminal task state, and \(\mathcal Q_D\) reads out the count. This notation
supports two competing hypotheses:
\paragraph{\(H_{\rm run}\): record-wise update.}
A predicate-conditioned record state is updated by approximately \(+1\) after
a valid occurrence and by \(0\) after a matched non-target, and the updated
state affects later computation.

\paragraph{\(H_{\rm agg}\): terminal aggregation.}
Local states may contain content or ordinal information, but the
task-relevant total is constructed mainly near \texttt{Total:} through broad
multi-record access and terminal readout.
The alternatives are not defined by whether a probe decodes \(N\). They are
distinguished by when a controlled \(+1\) becomes available and where a
natural activation patch transports it.

For CoT or explicit enumeration,
\begin{equation}
\begin{aligned}
  \pi_k &= \Pi_C(q,\Gamma_{k-1}),&
  e_k &= \mathcal W_C(R_{\pi_k}),\\
  \Gamma_k &= \mathcal F_C(\Gamma_{k-1},e_k),&
  \widehat N_C &= \mathcal Q_C(\Gamma_K).
\end{aligned}
\label{eq:cot-program}
\end{equation}
\(\pi_k\) is the source occurrence actually selected at step \(k\), \(e_k\)
is the written evidence, and \(\Gamma_k\) is the progress carrier used to
choose another source or stop. \(\Gamma_k\) may combine residual state
\(P_k\), the visible trace, and its KV cache; it is not assumed to be a scalar
neural counter. The registered semantic label is \(p_k\) from
\cref{eq:semantic-progress}, not the visible index \(k\).

\subsection{When we use the word counter}

A registered state is called a \emph{counter} only after three nested tests:
\begin{enumerate}
  \item \textbf{Decodable.} A decoder trained on disjoint families predicts
  the registered cumulative count, semantic progress, or final total beyond
  nuisance-only and token/position controls.
  \item \textbf{Dynamically consistent.} Direct record states exhibit the
  sibling-controlled \(+1/0\) profile; CoT progress increases after a newly
  retrieved valid occurrence and not after a matched non-target or repeated
  occurrence.
  \item \textbf{Causally used.} A single-donor natural patch or a
  validation-selected steering intervention moves a registered downstream
  endpoint in the predicted direction.
\end{enumerate}
A state that passes only the first test is a count-decodable
representation. PCA is a visualization and satisfies none of these criteria
by itself. Terminal states \(Z_D\) and \(Z_C\) are not called running counters;
they require total decodability and causal transport to the complete answer.

\subsection{Held-out decoding and geometry}

The token landmarks and family-equal analysis unit are fixed in
\cref{tab:landmarks,eq:family-weighting}. At each layer and semantic role, a
nuisance-only model \(g_0(u)\) is compared with a state-augmented model
\(g_1(h^\ell,u)\). For Direct local states,
\[
  u_D=(i,K,N,L,a_i/L,N/L,\text{depth bin},\text{template ID});
\]
for CoT progress,
\[
  u_C=(k,N,L,\text{trace-token length},\text{visible-label identity}).
\]
Remaining count \(N-c_q(a_i)\) is a separate outcome, not placed in the same
nuisance vector as both \(N\) and \(c_q(a_i)\). Evidence for a state is the
held-out improvement of \(g_1\) over \(g_0\), not raw predictability of a
count correlated with position or density.

For visualization, one shared PCA basis per model and layer is fitted on a
discovery sample balanced over registered role, interface, semantic label,
and \((N,L)\) cell; the downsampling seed is stored. The frozen basis projects
individual locked-test states. Mode-specific PCA is a sensitivity analysis.
Correct and incorrect outputs are both included in the primary fit and may be
stratified only for a registered diagnostic. No trajectory is created by
averaging states within a prompt or by selecting successful traces.

\begin{figure}[H]
\centering
\fbox{\parbox[c][5.0cm][c]{0.94\textwidth}{
\centering
\textbf{TODO Figure: Registered states and causal transport}\\[0.5em]
\begin{minipage}[t]{.30\linewidth}\centering
\textit{Held-out geometry}\\
individual \(H_i,P_k,Z_D,Z_C\)\\
shared PCA basis\\
family centroids as summaries
\end{minipage}\hfill
\begin{minipage}[t]{.30\linewidth}\centering
\textit{Dynamic tests}\\
Direct prefix-shift \(0/1/0\)\\
CoT new/invalid/repeated\\
predicate-after control
\end{minipage}\hfill
\begin{minipage}[t]{.30\linewidth}\centering
\textit{Causal transport}\\
single-donor patching\\
paired geometric steering\\
complete sequence endpoint
\end{minipage}\\[0.6em]
\small Qwen3-8B and Gemma4-E4B use the same registered variables but
independently selected layers. Cross-role dispersion is descriptive only.
}}
\caption{Planned representation figure. Readability, correct dynamics, and
causal transport are shown separately; only their conjunction licenses a
counter claim.}
\label{fig:counter-todo}
\end{figure}

\begin{todobox}
\textbf{R1---registered-state decoding.}
For every state in \cref{tab:landmarks}, fit family-weighted nuisance-only and
state-augmented decoders for Qwen3-8B and Gemma4-E4B. Freeze preprocessing,
ridge penalties, PCA bases, and layer selection before the locked test.
Report incremental held-out \(R^2\)/MAE, unseen identity and template
generalization, \((N,L)\) transfer, conditional dispersion, and correct/error
strata. Plot individual locked-test points plus family centroids; do not fit a
projection to the presentation prompts.
\end{todobox}

\subsection{Direct: running update or terminal aggregation?}

For predicate-before and predicate-repeat prompts, extract \(H_i^\ell\) at all
\(K\) controlled slots, not only positives. The primary dynamic test uses
prefix-shift siblings \(A\) and \(B\) that exchange a positive at \(r\) and a
matched negative at \(s>r\). Let \(\widetilde H_j^{A,\ell}\) and
\(\widetilde H_j^{B,\ell}\) be nuisance-residualized states at the same
unchanged landmark, and let \(g_h^\ell\) be the hidden-only decoder frozen on
discovery families. Define
\begin{equation}
  \delta_j^\ell=
  g_h^\ell(\widetilde H_j^{A,\ell})
  -
  g_h^\ell(\widetilde H_j^{B,\ell}).
\label{eq:prefix-shift}
\end{equation}
\(H_{\rm run}\) predicts
\[
  \delta_j^\ell\simeq
  \begin{cases}
    0,&j<r,\\
    1,&r\leq j<s,\\
    0,&j\geq s.
  \end{cases}
\]
This same-position contrast fixes final \(N\), record identity, local token,
length, and location. Adjacent-position increments inside one prompt are
secondary content/update diagnostics. Predicate-after states provide the
causal-mask negative control.

For causal transport, a total-shift pair differs at one equal-length slot
\(r\), but the primary patch is applied at the first \emph{unchanged}
downstream delimiter \(j>r\). Its local token and suffix are identical across
siblings, reducing record-content transport. Patching the modified record's
own endpoint is retained only as a positive evidence-transport control.
Endpoints are decoded cumulative count at later landmarks, \(Z_D\), complete
continuation likelihood after \texttt{Total:}, and free-generated exact
count/MAE. Identity-preserving positive-to-positive patches, matched
negative-to-negative patches, unrelated positions, random states,
norm-matched states, and sham patches are controls.

\(H_{\rm agg}\) also requires positive evidence; it is not accepted merely
because the local test fails. We test whether query-local multi-record routes
carry \(N\) into \(Z_D\), whether total-shift terminal patches move the
complete answer, and whether those effects survive controls for local ordinal
information. Evidence for \(H_{\rm run}\) requires decodability, the
\(0/1/0\) dynamic profile, and propagated downstream patches. Evidence for
\(H_{\rm agg}\) requires a selectively causal terminal route/state
intervention.

\begin{todobox}
\textbf{R2---Direct dynamics and patches.}
At every frozen candidate layer, run prefix-shift and total-shift experiments
on locked families. Use the same-position hidden-only contrast at unchanged
landmarks for the primary \(+1/0\) test; do not treat a second independent
valid occurrence as a no-op. Patch single natural donors at unchanged
downstream \(H_j\) and at \(Z_D\) separately. Measure later cumulative-state
decoding, terminal state, complete-answer likelihood, and free-generation
exact count/MAE. Include predicate-after and predicate-repeat,
identity-preserving, position, random, norm, and sham controls. Report
separate evidence for \(H_{\rm run}\) and \(H_{\rm agg}\); a failed local
counter test is not by itself evidence for aggregation.
\end{todobox}

\subsection{CoT progress without label leakage}

At step \(k\), let the registered next source in passage order be
\begin{equation}
  g_k=
  \operatorname*{arg\,min}_{i\in G_q\setminus\mathcal V_{k-1}} s_i,
\label{eq:gold-next}
\end{equation}
where \(s_i\) is source position. If the set is empty, the target is
\textsc{stop}. Gold-target alignment is the primary estimand; attention to
the source the model actually emits, \(\pi_k\), is reported separately as
execution alignment so a consistently wrong retrieval cannot receive a high
mechanism score.

For each visible step \(k\), we teacher-force three equal-length prefixes whose
newest item is a new gold source (\(\Delta p=1\)), a matched non-target
(\(\Delta p=0\)), or a repeat of an already visited occurrence
(\(\Delta p=0\)). Complementary prefixes hold \(p_k\) fixed while changing
trace extent. Numeric, random nonnumeric, permuted, and absent labels appear
outside the invariant \texttt{Item:} and \(D_{\rm eor}\) tokens. A state that
tracks \(p_k\) under these transformations is evidence for semantic progress;
a state that tracks only \(k\) or the displayed label is scaffolding.

A clean \(P_k^\ell\) is patched into a matched corrupt prefix before the model
chooses its next source. Proximal endpoints are the probability and generated
identity of \(g_{k+1}\), duplicate probability, and continue/stop probability;
later endpoints are unique coverage, \(P_{k+1}\), terminal state, and the
complete count. Residual-only transport patches \(P_k^\ell\) while preserving
the receiver trace/KV cache. Trace/KV-only path patching transports donor
trace keys and values while preserving the receiver residual. These carriers
are complementary and are not conflated into one scalar counter.

For geometric steering, the discovery direction is the normalized paired
difference
\begin{equation}
  v_\ell=
  \frac{\E[P_k^\ell\mid\Delta p=1]-\E[P_k^\ell\mid\Delta p=0]}
       {\|\E[P_k^\ell\mid\Delta p=1]-\E[P_k^\ell\mid\Delta p=0]\|_2}.
\label{eq:progress-steering}
\end{equation}
The layer and dose are selected on validation families. Norm-matched random
and orthogonal directions are controls. Steering defines one direction from a
cross-prompt mean; primary patching still uses a single matched natural donor.

\begin{todobox}
\textbf{R3---CoT anti-leakage and state control.}
Run the equal-length new-gold/non-target/repeated-source prefix experiment
with unindexed, random-label, and permuted-index interfaces. Test same-\(k\)
different-\(p_k\) and same-\(p_k\) different-extent pairs. Separately
transport residual state, trace/KV state, and both; run single-donor natural
patches and a validation-selected steering sweep. Report gold next-source
likelihood and identity, emitted-source alignment, duplicate and stop
behavior, later progress, unique coverage, and the complete parsed count.
Require effects beyond norm-matched, random, orthogonal, position, and sham
controls before calling \(P_k\) a semantic progress carrier.
\end{todobox}

\subsection{Conditional state dispersion is not probe error}

Let \(\widetilde h^\ell\) be the residual after a discovery-set nuisance model
removes registered position, length, density, template, and label covariates.
For a semantic count or progress value \(c\), define family-equally weighted
moments
\[
  \mu_c^\ell=\E[\widetilde h^\ell\mid c],
  \qquad
  \Sigma_c^\ell=\Cov(\widetilde h^\ell\mid c).
\]
For adjacent values, let
\[
  d_c=\mu_{c+1}^\ell-\mu_c^\ell,\qquad
  u_c=\frac{d_c}{\|d_c\|_2+\epsilon},\qquad
  \overline\Sigma_c=\frac{\Sigma_c^\ell+\Sigma_{c+1}^\ell}{2}.
\]
The registered local confusion statistic is
\begin{equation}
  \nu_{\rm state}^\ell(c)=
  \frac{u_c^\top\overline\Sigma_c u_c}
       {\|d_c\|_2^2+\epsilon}.
\label{eq:state-snr}
\end{equation}
It is reported only when adjacent separation exceeds a validation-frozen
threshold. Held-out decoder error measures accessibility to a chosen decoder
and is reported separately. We test whether Direct terminal dispersion or
error grows with \(N,L\), and whether controlled CoT terminal states are less
confusable. Because \(H_i\), \(P_k\), and \(Z\) occupy different token roles,
cross-role dispersion comparisons are exploratory and cannot establish an
algorithmic path change.
